\chapter{ETF, contrats à terme et arbitrage d'indice}\label{ch:iv2:etfs-futures-and-index-arbitrage}

\begin{ivtrains}
L'arithmétique du desk delta-one~: à quel prix un ETF devrait se traiter, quand la création ou le rachat devient rentable, ce que contient le prix d'un contrat à terme, ce que fait le recalibrage quotidien d'un fonds à effet de levier sur lui-même et sur le marché, et ce que valent un rééquilibrage ou un roll. Sauf indication contraire dans l'énoncé, les taux sont ici simples.
\end{ivtrains}

\begin{question}[$\star$ \iqroles{QT, Trader} \iqfirm{teneur de marché, hedge fund, banque} \iqlevel{junior}]\label{iq:iv2:etfs-futures-and-index-arbitrage:1}
Pour chaque lot de 100 parts de l'ETF, celui-ci détient 50 actions d'une valeur cotée à 20~\$, 30 actions d'une valeur cotée à 50~\$, 10 actions d'une valeur cotée à 100~\$, et 500~\$ de liquidités. Quelle est sa valeur liquidative par part~? L'ETF se traite à 40.10~\$~: quelle est sa prime en points de base~?
\end{question}

\begin{question}[$\star$ \iqroles{QT, Trader} \iqfirm{teneur de marché, hedge fund} \iqlevel{junior}]\label{iq:iv2:etfs-futures-and-index-arbitrage:2}
La valeur liquidative d'un ETF est de 40.00~\$. Acheter le panier et procéder à la création coûte au total 6 cents par part d'ETF, et racheter et vendre le panier coûte au total 4 cents. Entre quels cours de l'ETF n'y a-t-il pas d'arbitrage~? Le cours acheteur est de 40.09~\$~: que faites-vous, et combien gagnez-vous par part~?
\end{question}

\begin{question}[$\star$ \iqroles{QT, Trader, QR} \iqfirm{banque, hedge fund, teneur de marché} \iqlevel{junior}]\label{iq:iv2:etfs-futures-and-index-arbitrage:3}
Une action se traite à 100~\$. Elle verse un dividende de 1~\$ dans trois mois. Le taux d'intérêt est de 4~\% par an, et vous empruntez et prêtez à ce taux. Quel est le prix juste d'un forward à six mois~? Indiquez les transactions qui permettent de verrouiller un profit si le forward se traite à 102~\$.
\end{question}

\begin{question}[$\star$ \iqroles{QT, Trader, Risque} \iqfirm{hedge fund, banque} \iqlevel{junior}]\label{iq:iv2:etfs-futures-and-index-arbitrage:4}
Un indice vaut 4\,000, et son contrat à terme a un multiplicateur de 50~\$ par point. Quelle est la valeur notionnelle d'un contrat~? Combien de contrats couvrent un portefeuille de 10~millions de dollars avec un bêta de 1, et combien avec un bêta de 1.2~?
\end{question}

\begin{question}[$\star$ \iqroles{QT, Trader} \iqfirm{banque, hedge fund, teneur de marché} \iqlevel{junior}]\label{iq:iv2:etfs-futures-and-index-arbitrage:5}
Un indice vaut 4\,000. Son contrat à terme, à échéance dans un trimestre, se traite à 4\,020. Le rendement du dividende est de 1~\% par an. Quel taux de financement ce contrat à terme implique-t-il~? Si vous pouvez emprunter à 2.5~\%, que feriez-vous~?
\end{question}

\begin{question}[$\star$ \iqroles{QT, Trader} \iqfirm{banque, hedge fund, teneur de marché} \iqlevel{junior}]\label{iq:iv2:etfs-futures-and-index-arbitrage:6}
Un indice vaut 4\,000, et son contrat à terme à trois mois se traite à 4\,030. Le taux d'intérêt est de 5~\% par an. Quels dividendes, en points d'indice et sous forme de rendement annuel, le contrat à terme implique-t-il~?
\end{question}

\begin{question}[$\star$ \iqroles{QT, Trader, QR, Analyste} \iqfirm{toutes} \iqlevel{junior}]\label{iq:iv2:etfs-futures-and-index-arbitrage:7}
Un indice monte de 2~\%, puis baisse de 2~\%. Quel est son rendement total sur les deux jours~? Un fonds se recalibre chaque jour sur trois fois la performance quotidienne de l'indice. Quelle est sa performance~? Comment se compare-t-elle à trois fois le rendement de l'indice~?
\end{question}

\begin{question}[$\star$ \iqroles{QT, Trader, Analyste} \iqfirm{toutes} \iqlevel{junior}]\label{iq:iv2:etfs-futures-and-index-arbitrage:8}
Un indice passe de 100 à 110, puis revient à 100. Un fonds se recalibre chaque jour sur moins une fois la performance quotidienne de l'indice. Quel est le rendement total du fonds~? Quelqu'un le détient pour couvrir une position acheteuse sur l'indice pendant ces deux jours. Quel est le résultat combiné~?
\end{question}

\begin{question}[$\star$ \iqroles{QT, QR, Analyste} \iqfirm{hedge fund, banque} \iqlevel{junior}]\label{iq:iv2:etfs-futures-and-index-arbitrage:9}
Un fonds détient deux actions à parts égales, 100 millions de dollars dans chacune. L'une monte de 10~\%, l'autre baisse de 10~\%. Quelles transactions le ramènent à des pondérations égales, et quel est le taux de rotation à sens unique, rapporté au fonds~? Quelles transactions un fonds pondéré par la capitalisation devrait-il effectuer~? Pourquoi~?
\end{question}

\begin{question}[$\star\star$ \iqroles{QT, Trader, Risque} \iqfirm{hedge fund, banque} \iqlevel{junior}]\label{iq:iv2:etfs-futures-and-index-arbitrage:10}
Vous êtes long de 10 contrats à terme sur indice à 4\,000, avec un multiplicateur de 50~\$, et vous déposez 5~\% du notionnel en marge initiale. L'indice baisse à 3\,900 à la clôture. Quelle est votre perte, et quel montant de liquidités change de mains aujourd'hui~? Quelle est votre marge si la marge initiale est toujours exigée au nouveau prix~? Quelle est la différence entre cette position et la même exposition au moyen d'un forward~?
\end{question}

\begin{question}[$\star\star$ \iqroles{QT, Trader, Analyste} \iqfirm{toutes} \iqlevel{junior}]\label{iq:iv2:etfs-futures-and-index-arbitrage:11}
Un fonds reste exposé à la hausse en détenant le contrat à terme mensuel sur le pétrole le plus proche et en faisant le roll chaque mois. Le contrat le plus proche se traite à 70, et le suivant à 71.40. Le spot et l'ensemble de la courbe restent exactement inchangés pendant un an. Quel gain le fonds tire-t-il du roll au bout d'un mois~? Et sur un an~? Pourquoi un prix du pétrole inchangé ne se traduit-il pas par un résultat inchangé~?
\end{question}

\begin{question}[$\star\star$ \iqroles{QT, Trader, QR} \iqfirm{teneur de marché, hedge fund} \iqlevel{junior, confirmé}]\label{iq:iv2:etfs-futures-and-index-arbitrage:12}
Un fonds disposant de 1 milliard de dollars d'actifs se recalibre chaque jour sur trois fois l'indice. L'indice monte de 1~\% aujourd'hui. Quels sont ses actifs et son exposition avant le recalibrage, et combien doit-il acheter ou vendre à la clôture~? Écrivez une formule pour un levier $L$, une taille de fonds $A$ et un rendement $r$. Que fait un fonds à moins trois le même jour~?
\end{question}

\begin{question}[$\star\star$ \iqroles{QR, QT} \iqfirm{hedge fund, teneur de marché} \iqlevel{junior, confirmé}]\label{iq:iv2:etfs-futures-and-index-arbitrage:13}
Un fonds se recalibre quotidiennement sur $L$ fois un indice dont la volatilité annuelle est $\sigma$, sans taux ni frais. Sur un an, comment son rendement se rapporte-t-il à celui de l'indice~? Si l'indice a une volatilité de 30~\% et termine l'année sans variation, quel est le rendement du fonds à trois fois et celui du fonds à moins une fois~?
\end{question}

\begin{question}[$\star\star$ \iqroles{QT, QR, Dev} \iqfirm{banque, hedge fund} \iqlevel{junior}]\label{iq:iv2:etfs-futures-and-index-arbitrage:14}
Un indice pondéré par les cours compte trois actions à 100~\$, 50~\$ et 30~\$, avec un diviseur de 3, si bien qu'il vaut 60. La première action subit une division d'actions de 2 pour 1. Quel doit être le nouveau diviseur, et pourquoi~? Si les trois actions progressent ensuite de 10~\%, quelle est la valeur de l'indice, et quelle aurait-elle été sans la division~?
\end{question}

\begin{question}[$\star\star$ \iqroles{QT, Trader} \iqfirm{banque, hedge fund, teneur de marché} \iqlevel{junior}]\label{iq:iv2:etfs-futures-and-index-arbitrage:15}
Un indice de prix vaut 4\,000. Une action dont le poids est de 2~\% passe ex-dividende aujourd'hui, en versant 1~\% de son cours. De combien de points l'indice baisse-t-il à l'ouverture, toutes choses égales par ailleurs~? Que se passe-t-il pour l'indice de rendement total, et pour la fair value du contrat à terme échéant après cette date~?
\end{question}

\begin{question}[$\star\star$ \iqroles{QT, Trader} \iqfirm{banque, hedge fund} \iqlevel{junior, confirmé}]\label{iq:iv2:etfs-futures-and-index-arbitrage:16}
Un indice vaut 4\,000. Vous achetez le panier et vendez un contrat à terme à un an, à un prix qui suppose des dividendes de 2~\% du spot. Les dividendes s'élèvent finalement à 2.5~\%. Quel est votre gain en points~? Qui est de l'autre côté de ce risque~? Que négocieriez-vous si vous ne vouliez que cette exposition aux dividendes~?
\end{question}

\begin{question}[$\star\star$ \iqroles{QT, Trader, Analyste} \iqfirm{hedge fund, banque} \iqlevel{junior, confirmé}]\label{iq:iv2:etfs-futures-and-index-arbitrage:17}
Vous souhaitez une exposition à l'indice. Un contrat à terme que l'on fait rouler coûte 15 points de base par an de financement implicite au-dessus de votre taux de financement, plus 4 points de base par an de coûts de roll. Un ETF facture 9 pb de frais annuels, rapporte 2 pb par an grâce au prêt de titres et coûte 1 point de base de fourchette à l'achat comme à la vente. Sur un an, lequel est le moins cher~? Pour quelle durée de détention sont-ils égaux~?
\end{question}

\begin{question}[$\star\star$ \iqroles{QT, Trader} \iqfirm{banque, hedge fund, teneur de marché} \iqlevel{junior, confirmé}]\label{iq:iv2:etfs-futures-and-index-arbitrage:18}
Un contrat à terme sur indice échéant dans trois mois se traite à 4\,020, et le suivant, trois mois plus tard, à 4\,050. Les taux sont à 5~\% par an, et les dividendes sur ce dernier trimestre s'établissent à 2~\% par an. Le calendar spread est-il juste, en prenant 4\,020 comme base~? Et si le contrat le plus éloigné se traite à 4\,060~? Que faites-vous si vous êtes long du contrat le plus proche et devez le rester~?
\end{question}

\begin{question}[$\star\star$ \iqroles{Analyste, QT} \iqfirm{hedge fund, banque} \iqlevel{junior}]\label{iq:iv2:etfs-futures-and-index-arbitrage:19}
Un fonds indiciel facture des frais de 0.20~\% par an et gagne 0.08~\% grâce au prêt de titres. Il conserve 1~\% de ses actifs en liquidités, et sur l'année l'indice rapporte 8~\% tandis que les liquidités rapportent 4~\%. Il détient un échantillon de l'indice dont le rendement s'en écarte de 5 points de base par an, avec une espérance nulle. Quel est son écart de rendement attendu avec l'indice~? Distinguez le tracking difference de la tracking error.
\end{question}

\begin{question}[$\star\star$ \iqroles{QT, Trader} \iqfirm{banque, hedge fund} \iqlevel{junior, confirmé}]\label{iq:iv2:etfs-futures-and-index-arbitrage:20}
Un indice vaut 4\,000, avec un contrat à terme à six mois et des dividendes de 2~\% par an. Vous empruntez à 5~\% et prêtez à 4.5~\%. La vente à découvert du panier coûte en outre des frais de 0.3~\% par an. L'entrée et la sortie de la transaction coûtent 0.1~\% du spot au total. Entre quels cours à terme n'y a-t-il pas d'arbitrage pour vous~? Pourquoi la bande n'est-elle pas symétrique~?
\end{question}

\begin{question}[$\star\star$ \iqroles{QT, Trader, Risque} \iqfirm{teneur de marché, hedge fund, banque} \iqlevel{junior, confirmé}]\label{iq:iv2:etfs-futures-and-index-arbitrage:21}
Lors d'un mouvement de vente massif sur le crédit, un ETF obligataire se traite à 95 alors que sa valeur liquidative publiée est de 100. Vous recevriez des obligations qui, vendues immédiatement, rapporteraient réellement 95.50, et le rachat coûte 0.20. Que gagnez-vous par part d'ETF si tout se produit en même temps~? La vente des obligations prend trois jours, avec une volatilité quotidienne de 0.5~\% de leur cours. S'agit-il d'un arbitrage~? Lequel des deux chiffres, 95 ou 100, est la meilleure estimation de la valeur~?
\end{question}

\begin{question}[$\star\star$ \iqroles{QT, Trader} \iqfirm{teneur de marché} \iqlevel{junior, confirmé}]\label{iq:iv2:etfs-futures-and-index-arbitrage:22}
Vous cotez un ETF à 40.00~\$ et 40.02~\$, la fair value étant de 40.01~\$. Une fois exécuté, vous couvrez aussitôt avec un contrat à terme, en payant la moitié de sa fourchette, soit 0.5 point de base. Si votre couverture et votre ETF s'écartent pendant que vous les détenez, quel est le risque~? Quel est votre résultat net par action sur une exécution~? Pour 100\,000 actions, combien cela représente-t-il en dollars, et que coûte un mouvement de base d'un point de base défavorable~?
\end{question}

\begin{question}[$\star\star$ \iqroles{QT, Trader, Analyste} \iqfirm{teneur de marché, hedge fund} \iqlevel{junior, confirmé}]\label{iq:iv2:etfs-futures-and-index-arbitrage:23}
Acheter le panier d'un ETF coûte une demi-fourchette de 2 points de base. Sur une journée, les investisseurs achètent 100 millions de dollars de l'ETF et en vendent 90 millions, et un teneur de marché prend la contrepartie de l'ensemble. Que doit couvrir le teneur de marché dans le panier~? Combien cela coûte-t-il par dollar d'ETF échangé~? Pourquoi la fourchette d'un ETF peut-elle être plus étroite que la fourchette des titres qu'il détient~?
\end{question}

\begin{question}[$\star\star$ \iqroles{QT, Trader, Analyste} \iqfirm{banque, hedge fund} \iqlevel{junior, confirmé}]\label{iq:iv2:etfs-futures-and-index-arbitrage:24}
Vous souhaitez une exposition à une action dont le rendement est de 2~\% par an. La détenir au comptant vous coûte le financement plus 30 points de base, et vous perdez 15~\% du dividende au titre de la retenue à la source. La détenir par un total return swap coûte le financement plus 40 points de base, et le dealer reverse 95~\% du dividende. Lequel est le moins cher, et de combien~? Quel risque prenez-vous avec le swap, et que fait le dealer pour se couvrir~?
\end{question}

\begin{question}[$\star\star\star$ \iqroles{QT, Trader, Analyste} \iqfirm{hedge fund, teneur de marché} \iqlevel{junior, confirmé}]\label{iq:iv2:etfs-futures-and-index-arbitrage:25}
Un titre dont la valeur de marché est de 10 milliards de dollars, dont 80~\% constituent le flottant, entre dans un indice pondéré par la capitalisation. Les fonds qui répliquent l'indice détiennent ensemble 20~\% du flottant de chaque composante. Que doivent-ils acheter~? Le volume quotidien est de 100 millions de dollars~: combien de jours cela représente-t-il à 100~\% du volume, et à 20~\%~? Avec la règle selon laquelle le coût est $\sigma\sqrt{Q/V}$ et une volatilité quotidienne de 2~\%, que fait le cours, et pourquoi devriez-vous douter de ce chiffre~?
\end{question}

\begin{question}[$\star\star\star$ \iqroles{QR, QT} \iqfirm{hedge fund, teneur de marché} \iqlevel{junior, confirmé}]\label{iq:iv2:etfs-futures-and-index-arbitrage:26}
Un collègue remarque qu'un fonds à trois fois et un fonds à moins trois fois perdent tous deux de la valeur sur les marchés volatils, et propose de vendre à découvert 1~\$ de chacun et de conserver la position pendant un an, en ignorant les coûts d'emprunt. Les rendements quotidiens sont d'espérance nulle et la volatilité annuelle est de 30~\%. Quel est votre profit espéré~? Quel est votre profit médian, et à quelle fréquence perdez-vous de l'argent~? Qu'est-ce que cela dit du «~decay~»~?
\end{question}

\begin{question}[$\star\star\star$ \iqroles{QT, QR, Trader} \iqfirm{teneur de marché, hedge fund} \iqlevel{junior, confirmé}]\label{iq:iv2:etfs-futures-and-index-arbitrage:27}
Les fonds à trois fois détiennent 10~milliards de dollars sur un indice, et les fonds à moins trois fois en détiennent 5~milliards. L'indice recule de 4~\% aujourd'hui. Quel montant chaque groupe doit-il échanger à la clôture, et dans quel sens~? Vérifiez chacun en calculant les expositions. Si la liquidité du marché sur les 15 dernières minutes atteint 20~milliards de dollars, quelle part cela représente-t-il~? Pourquoi cela ressemble-t-il à la couverture d'une position short gamma~?
\end{question}

\begin{question}[$\star\star\star$ \iqroles{QR, QT} \iqfirm{banque, hedge fund} \iqlevel{confirmé}]\label{iq:iv2:etfs-futures-and-index-arbitrage:28}
Une action ne verse aucun dividende et vaut 100. Le taux à une période est de 4~\%. Sous la probabilité risque-neutre, à la période suivante elle vaut soit 120 avec un taux à une période de 6~\%, soit 84 avec un taux à une période de 2~\%, et, dans chaque état, elle croît ensuite à ce taux pendant une période supplémentaire. Déterminez les probabilités risque-neutres. Quels sont le prix forward et le prix du contrat à terme à deux périodes~? Pourquoi sont-ils différents, et qu'est-ce qui annule cette différence~?
\end{question}

\begin{question}[$\star\star\star$ \iqroles{QR, QT, Risque} \iqfirm{hedge fund, banque} \iqlevel{junior, confirmé}]\label{iq:iv2:etfs-futures-and-index-arbitrage:29}
Un ETF sectoriel présente une volatilité annuelle de 18~\%. Vous le couvrez avec deux actions, $A$ de volatilité 20~\% et $B$ de volatilité 30~\%, de corrélation 0.5 entre elles. L'ETF a une corrélation de 0.8 avec $A$ et de 0.7 avec $B$. Quelles pondérations de couverture de variance minimale utilisez-vous par dollar d'ETF, et quelle volatilité résiduelle subsiste-t-il~? Pourquoi ne pas couvrir avec 100 valeurs~?
\end{question}

\begin{question}[$\star\star\star$ \iqroles{QR, QT, Trader} \iqfirm{teneur de marché, hedge fund} \iqlevel{junior, confirmé}]\label{iq:iv2:etfs-futures-and-index-arbitrage:30}
Un future est trop cher de 6 points par rapport à sa fair value, et il doit converger vers elle à l'échéance, dans trois jours. L'écart monte ou descend de 4 points chaque jour, avec des probabilités égales les jours 1 et 2, puis devient nul le jour 3. Vous pouvez détenir une position de $+1$, $0$ ou $-1$ sur l'écart. Si vous restez short d'aujourd'hui jusqu'à l'échéance, quel gain réalisez-vous~? Quelle est l'espérance maximale de votre gain avec des retournements sans coût~? Combien vaut cette option si chaque retournement coûte 0.5 point au total~? Quand l'utilisez-vous~?
\end{question}

\begin{question}[$\star\star\star$ \iqroles{QT, Trader, Risque} \iqfirm{hedge fund, teneur de marché} \iqlevel{junior, confirmé}]\label{iq:iv2:etfs-futures-and-index-arbitrage:31}
Un fonds sur matières premières cesse d'émettre de nouvelles actions parce qu'il a atteint une limite sur ses avoirs, tandis que les anciennes actions continuent de se négocier. La demande est forte, et les actions atteignent une prime de 20~\% par rapport à la valeur liquidative. Vous achetez à ce niveau. Un mois plus tard, le sous-jacent progresse de 5~\%, la création reprend et la prime disparaît. Quel rendement dégagez-vous~? Pourquoi l'arbitrage habituel ne referme-t-il pas la prime~? Que feriez-vous avec une position short~?
\end{question}

\begin{question}[$\star\star\star$ \iqroles{QT, Trader, Risque, Analyste} \iqfirm{hedge fund, banque} \iqlevel{junior, confirmé}]\label{iq:iv2:etfs-futures-and-index-arbitrage:32}
Un ETF a 100 millions d'actions en circulation et un encours vendeur à découvert de 120 millions d'actions. Comment peut-on vendre à découvert plus de 100~\%~? Combien d'actions les porteurs inscrits et ceux qui sont longs par le biais de prêts croient-ils détenir au total~? Que font les teneurs de marché lorsqu'ils livrent~? Quel est le risque réel pour les investisseurs si l'ETF est liquidé~?
\end{question}

\begin{question}[$\star\star\star$ \iqroles{Analyste, QT} \iqfirm{hedge fund, banque} \iqlevel{junior, confirmé}]\label{iq:iv2:etfs-futures-and-index-arbitrage:33}
Un fonds réplique un indice pondéré par la capitalisation. Une action entre dans l'indice avec un poids de 0.1~\%. Un achat à la clôture effective devrait entraîner un impact de 3~\% sur cette action. Un achat réparti sur les cinq jours précédents coûte 1~\%, mais le fonds détient ensuite une position supplémentaire, que vous devez prendre égale à la totalité des 0.1~\% pendant les cinq jours. L'action a une volatilité quotidienne de 2~\%, couverte uniquement par le reste de l'indice. Quelle est l'économie espérée pour l'ensemble du fonds, en points de base~? Quel risque de tracking error le fonds prend-il~? Laquelle choisiriez-vous~?
\end{question}

\begin{question}[$\star\star\star\star$ \iqroles{QR} \iqfirm{hedge fund, teneur de marché, banque} \iqlevel{confirmé}]\label{iq:iv2:etfs-futures-and-index-arbitrage:34}
Un fonds vise $L$ fois le rendement d'un indice à volatilité constante $\sigma$ et taux $r$, rééquilibré en continu, avec des frais $f$. En utilisant le lemme d'Itô, dérivez sa valeur à l'instant $T$ en fonction de l'indice. Pour $L = 3$, $\sigma = 20~\%$, $T = 1$, sans frais et un fonds initialisé à 1, quelle est sa valeur si l'indice est inchangé et les taux nuls, et quelle est-elle si l'indice monte de 10~\% et que les taux sont de 4~\%~? Avec des taux nuls, combien vaut un call à un an à la monnaie sur le fonds, pour 100 de valeur du fonds~? Comparez-le à trois calls à un an à la monnaie sur un indice à 100.
\end{question}

\begin{question}[$\star\star\star\star$ \iqroles{QT, QR, Trader} \iqfirm{hedge fund} \iqlevel{confirmé}]\label{iq:iv2:etfs-futures-and-index-arbitrage:35}
Des fonds doivent acheter 1.6~milliard de dollars d'une action à la clôture dans cinq séances de bourse, à la suite de son annonce d'entrée dans un indice. Le volume quotidien est de 100~millions de dollars et la volatilité quotidienne de 2~\%. Utilisez une pression sur les prix au pic de $\sigma\sqrt{Q/V}$. La moitié de celle-ci a été intégrée dans le cours à l'annonce. Vous achetez 100~millions de dollars sur les cinq jours à 20~\% du volume, à un coût moyen de $\sigma\sqrt{Q/V}$ calculé sur votre propre taille, et vous vendez aux fonds à la clôture. Quel est votre profit espéré, et quel risque courez-vous~? Si d'autres acquièrent 400~millions de dollars d'ici la clôture et font de même, qu'advient-il de la pression~? Dans quels cas renonceriez-vous~?
\end{question}
